\documentclass[aspectratio=169,10pt]{beamer}

% ============================================================
% Packages
% ============================================================
\usepackage[utf8]{inputenc}
\usepackage[T1]{fontenc}
\usepackage{lmodern}
\usepackage{amsmath, amssymb}
\usepackage{booktabs}
\usepackage{array}
\usepackage{graphicx}
\usepackage{ragged2e}
\usepackage{tcolorbox}
% \tcbuselibrary{skins} % requires tikzfill (missing from TeX Live basic)
\usetikzlibrary{arrows.meta}
\usepackage[style=authoryear, backend=biber, maxcitenames=2, maxbibnames=99, uniquename=false, uniquelist=false]{biblatex}
\addbibresource{references.bib}
\usepackage{multicol}
\setbeamertemplate{bibliography item}{}
\ExecuteBibliographyOptions{giveninits=true, isbn=false, url=false, doi=false}
\DeclareNameAlias{default}{family-given}
\DeclareNameAlias{sortname}{family-given}
\renewcommand*{\bibfont}{\scriptsize}
\setlength{\bibitemsep}{2pt}
\renewbibmacro{in:}{}
\setbeamercolor{bibliography entry author}{fg=black}
\setbeamercolor{bibliography entry title}{fg=black}
\setbeamercolor{bibliography entry location}{fg=black}
\setbeamercolor{bibliography entry note}{fg=black}
% References not yet in references.bib are shown in red
\newcommand{\missingref}[1]{\textcolor{red}{#1}}

% ============================================================
% Colours (same as the written proposal)
% ============================================================
\definecolor{boxframe}{RGB}{150,58,78}
\definecolor{boxback}{RGB}{248,238,239}
\definecolor{darkgreen}{RGB}{0,100,0}
\definecolor{textgrey}{RGB}{110,110,110}

% ============================================================
% Beamer look
% ============================================================
\usetheme{default}
\setbeamertemplate{navigation symbols}{}
\setbeamersize{text margin left=0.6cm, text margin right=0.6cm}
\setbeamercolor{frametitle}{fg=boxframe}
\setbeamerfont{frametitle}{series=\bfseries, size=\Large}
\setbeamercolor{structure}{fg=boxframe}
\setbeamertemplate{itemize item}{\color{boxframe}--}
\setbeamertemplate{itemize subitem}{\color{boxframe}\textbullet}
\setbeamertemplate{enumerate item}{\color{boxframe}\insertenumlabel.}

% Header: frame title with a coloured rule underneath
\setbeamertemplate{frametitle}{%
  \vspace{0.4cm}%
  {\usebeamerfont{frametitle}\usebeamercolor[fg]{frametitle}\insertframetitle\par}
  \vspace{0.1cm}
  {\color{boxframe}\rule{\linewidth}{0.6pt}}
}

% Footer: course | title | page
\setbeamertemplate{footline}{%
  \hspace{0.4cm}{\scriptsize\color{textgrey}ENSAE -- IP Paris \;|\; Public Economics \;|\; \insertshorttitle}
  \hfill
  {\scriptsize\color{textgrey}\insertframenumber/\inserttotalframenumber}\hspace{0.4cm}
  \vspace{0.25cm}
}

% ============================================================
% Helpers
% ============================================================
% Grey placeholder text, to replace with content
\newcommand{\placeholder}[1]{{\color{textgrey}\itshape #1}}
% Coloured box, same style as the written proposal
\newtcolorbox{keybox}[1][]{%
  colback=boxback, colframe=boxframe, boxrule=0.6pt, arc=2pt,
  left=4pt, right=4pt, top=2pt, bottom=2pt,
  fonttitle=\bfseries\small, title={#1}
}

\newcommand{\blocklabel}[1]{\vspace{0.15cm}{\color{boxframe}\textbf{#1}}\par\vspace{0.05cm}}

% ============================================================
% Title information
% ============================================================
\title[\placeholder{Who pays when fathers don't?}]{\placeholder{Who pays when fathers don't?}}
\subtitle{\placeholder{Default intermediation of child support in France}}
\author{\placeholder{Eve Samani}, \placeholder{Isaure Pillet}, \placeholder{Antony Albergne}, \placeholder{Lucie Lang}}
\institute{Public Economics -- Enrico Rubolino \\ ENSAE -- IP Paris, 3A}
\date{Fall 2026}

% Title page layout
\setbeamertemplate{title page}{%
  \vfill
  \begin{center}
    {\color{boxframe}\rule{0.7\linewidth}{0.6pt}}\par
    \vspace{0.4cm}
    {\usebeamercolor[fg]{frametitle}\Huge\bfseries\inserttitle\par}
    \vspace{0.25cm}
    {\large\color{textgrey}\insertsubtitle\par}
    \vspace{0.4cm}
    {\color{boxframe}\rule{0.7\linewidth}{0.6pt}}\par
    \vspace{0.6cm}
    {\normalsize\insertauthor\par}
    \vspace{0.3cm}
    {\small\color{textgrey}\insertinstitute\par}
    \vspace{0.2cm}
    {\small\color{textgrey}\insertdate\par}
  \end{center}
  \vfill
}

\begin{document}

% ============================================================
% TITLE SLIDE (no footer, not counted in the slide numbers)
% ============================================================
\begin{frame}[plain, noframenumbering]
  \titlepage
\end{frame}



% ============================================================
% SLIDE 1 -- Motivation and Literature
% ============================================================

% ============================================================
% SLIDE 1 -- Motivation and Literature
% ============================================================

\begin{frame}{Motivation and Literature}

\vspace{-0.35cm}

\begin{columns}[T,onlytextwidth]
% ============================================================
% LEFT COLUMN -- MOTIVATION
% ============================================================
\begin{column}{0.48\textwidth}

  
\blocklabel{Motivation}
\footnotesize

\begin{itemize}
    \setlength{\itemsep}{0.09cm}
    \setlength{\topsep}{0pt}
    \justifying

    \item \textbf{Gendered costs of divorce:}
    Women's living standards fall by
    \textbf{\textcolor{boxframe}{19\%}}
    one year after divorce, compared with
    \textbf{\textcolor{boxframe}{2\%}}
    for men (Bonnet et al., 2021).

    \item \textbf{Unreliable child support:}
    Child support represents
    \textbf{\textcolor{boxframe}{18\%}}
    of single-parent families' resources,
    yet \textbf{\textcolor{boxframe}{30--40\%}}
    of payments were fully or partially unpaid.

    \item \textbf{Policy reform:}
    France progressively introduced
    \textbf{automatic public intermediation}
    of child support (March 2022;
    January 2023) to improve payment reliability.

\end{itemize}

\vspace{0.03cm}
{\tiny\color{textgrey}
Sources: Bonnet et al. (2021); CNAF (2023);
Ministry of Justice (2022).
}
\end{column}

% ============================================================
% RIGHT COLUMN -- RELATED LITERATURE
% ============================================================
\begin{column}{0.48\textwidth}

  \blocklabel{Related literature}
  \footnotesize

  \begin{itemize}
    \setlength{\itemsep}{0.09cm}
    \setlength{\topsep}{0pt}
    \justifying

    \item \textbf{Post-Separation Income \& Support}:
    \textcite{bonnet2021flip} 

    \item \textbf{Agency Failures \& Enforcement}:
    \textcite{weiss1993transfers};
    \textcite{chiappori2007divorce};
    \textcite{freeman2001dunning}.

    \item \textbf{Behavioral Responses \& IFPA}:
    \textcite{rossin2018parental};
    \textcite{cancian2013discouraging};
    \textcite{titlied2024} \textit{(French IFPA)}.

  \end{itemize}

\end{column}

\end{columns}


% ============================================================
% BOTTOM BOXES -- ALIGNED HORIZONTALLY
% ============================================================
\vspace{-0.25cm}

\begin{columns}[T,onlytextwidth]

% LEFT BOX -- RESEARCH QUESTION
\begin{column}{0.48\textwidth}
    \begin{keybox}[Research question]
        \footnotesize
        How does the shift to default public
        child-support intermediation affect
        payment compliance and mothers'
        economic security after separation?
    \end{keybox}
\end{column}

% RIGHT BOX -- OUR CONTRIBUTION
\begin{column}{0.48\textwidth}
    \begin{keybox}[Our contribution]
        \footnotesize
        France's phased reform provides a rare
        opportunity to investigate how public
        enforcement of private transfers shapes
        women's economic security after separation.
    \end{keybox}
\end{column}

\end{columns}
\end{frame}



% ============================================================
% SLIDE 2 -- Data
% ============================================================
\begin{frame}{Data}
\footnotesize
\begin{columns}[T,onlytextwidth]
  % ---------- Left column: source and sample ----------
  \begin{column}{0.44\textwidth}
    \blocklabel{Source}
    \begin{itemize}\setlength{\itemsep}{2pt}
      \item \textbf{ALLSTAT-FR6} (CNAF, accessed via CASD)
      \item Administrative panel, \textbf{2017--2025}
    \end{itemize}

    \blocklabel{Sample}
    \begin{itemize}\setlength{\itemsep}{2pt}
      \item Separated parents with $\geq 1$ child and child support due
      \item Separations within $\pm 12$ months of each cutoff (March 2022; January 2023)
    \end{itemize}

    \blocklabel{Limitation}
    \begin{itemize}\setlength{\itemsep}{2pt}
      \item No individual IFPA status or title date $\Rightarrow$ we identify the effect of
            \emph{exposure} to the default, not of take-up
    \end{itemize}
  \end{column}

  % ---------- Right column: variables ----------
  \begin{column}{0.53\textwidth}
    \blocklabel{Variables}
    \renewcommand{\arraystretch}{1.25}
    \begin{tabular}{@{}p{0.27\linewidth}p{0.66\linewidth}@{}}
      \toprule
      \textbf{Role} & \textbf{Definition} \\
      \midrule
      Running variable & Month of separation $c_i$ (\texttt{NBMSIFAM}) \\
      Group & Married vs.\ unmarried before separation (\texttt{SITFAM}) \\
      Outcomes $Y_i$ & Child support received; ASF; RSA; prime d'activit\'e; earnings \\
      Horizons & 6, 12 and 24 months after separation \\
      Covariates $X_i$ & Age, sex, number of children, pre-separation income, municipality \\
      \bottomrule
    \end{tabular}
  \end{column}
\end{columns}
\end{frame}

% ============================================================
% SLIDE 3 -- Methodology
% ============================================================
\begin{frame}{Methodology: Difference-in-Discontinuities}
\vspace{-0.2cm}
\footnotesize
\begin{columns}[T,onlytextwidth]
  % ---------- Left column: specification and interpretation ----------
  \begin{column}{0.49\textwidth}
    \begin{keybox}[{Specification, at each cutoff $c^*$}]
      \vspace{-0.25cm}
      \begin{align*}
      Y_i =\; & \underbrace{\textcolor{red}{\beta}\,\text{Post}_i \times \text{Treated}_i}_{\textcolor{red}{\text{effect of default IFPA}}}
      + \underbrace{\tau\,\text{Post}_i}_{\text{common jump}}
      + \underbrace{\theta\,\text{Treated}_i}_{\text{group gap}} \\
      & + \underbrace{f_g(c_i - c^*)}_{\text{cohort trends}}
      + \underbrace{X_i'\gamma}_{\text{controls}}
      + \varepsilon_i
      \end{align*}
      \vspace{-0.35cm}
    \end{keybox}

    \blocklabel{Causal evidence on default vs.\ opt-in enforcement}
    \justifying
    We estimate the effect of making public intermediation the \emph{default}, not of intermediation itself.
    $\textcolor{red}{\beta}$ is an ITT, which is policy-relevant because the default \emph{is} the reform.

    \blocklabel{Design}
    \justifying
    Two cutoffs where treated and control groups swap roles, which allows each design to be checked against the other.
  \end{column}

  % ---------- Right column: timeline, assumptions, checks ----------
  \begin{column}{0.49\textwidth}
    \centering
    \resizebox{\linewidth}{!}{%
    \begin{tikzpicture}[x=1.45cm, y=0.85cm, font=\small]
      \def\cutA{2.33}\def\cutB{4}\def\asf{3.67}\def\xend{6}
      \node[anchor=east, align=right] at (-0.2,2) {\textbf{Judicial divorces}\\[-2pt]{\scriptsize\color{textgrey}(married, divorced by a judge)}};
      \node[anchor=east, align=right] at (-0.2,0.8) {\textbf{Other separations}\\[-2pt]{\scriptsize\color{textgrey}(cohabiting, PACS, notarial divorces)}};
      \fill[textgrey!25] (0,1.75) rectangle (\cutA,2.25);
      \fill[boxframe!75] (\cutA,1.75) rectangle (\xend,2.25);
      \fill[textgrey!25] (0,0.55) rectangle (\cutB,1.05);
      \fill[boxframe!75] (\cutB,0.55) rectangle (\xend,1.05);
      \node[textgrey!80!black] at (\cutA/2,2) {\scriptsize opt-in};
      \node[white] at ({(\cutA+\xend)/2},2) {\scriptsize\bfseries default IFPA};
      \node[textgrey!80!black] at (\cutB/2,0.8) {\scriptsize opt-in};
      \node[white] at ({(\cutB+\xend)/2},0.8) {\scriptsize\bfseries default IFPA};
      \draw[-{Stealth}, thick] (0,0) -- (\xend+0.3,0) node[right] {\scriptsize cohort $c$};
      \foreach \x/\l in {0/2021, 2/2022, 4/2023, 6/2024}
        \draw (\x,0.08) -- (\x,-0.08) node[below] {\scriptsize \l};
      \draw[boxframe, very thick, dashed] (\cutA,-0.1) -- (\cutA,2.6) node[above] {\scriptsize\bfseries Mar.\ 2022};
      \draw[boxframe, very thick, dashed] (\cutB,-0.1) -- (\cutB,2.6) node[above] {\scriptsize\bfseries Jan.\ 2023};
      \draw[textgrey, thin, dotted] (\asf,-0.1) -- (\asf,1.4);
      \node[textgrey, anchor=north] at (\asf,-0.35) {\tiny ASF uprating};
      \node[draw=boxframe, fill=boxback, rounded corners, align=left, anchor=north, text width=5.2cm]
        at (0.9,-0.75) {\scriptsize \textbf{Cutoff 1}: treated = judicial divorces\\ control = other separations};
      \node[draw=boxframe, fill=boxback, rounded corners, align=left, anchor=north, text width=5.2cm]
        at (5.0,-0.75) {\scriptsize \textbf{Cutoff 2}: treated = other separations\\ control = judicial divorces (already treated)};
    \end{tikzpicture}}

    \raggedright
    \begin{itemize}\setlength{\itemsep}{1pt}
      \item \textbf{Assumptions}: (i) absent the reform, both groups jump equally at the cutoff; (ii) no sorting in separation timing or divorce type; (iii) no interaction with concurrent shocks
      \item \textbf{Checks}: no bunching in separations or jump in group shares; covariate balance; placebo cutoffs; staggered DiD
      \item \textbf{Caveat}: married $\neq$ judicial divorce $\Rightarrow$ conservative estimates
    \end{itemize}
  \end{column}
\end{columns}
\end{frame}

% ============================================================
% SLIDE 4 -- Expected effects: incidence of non-payment
% ============================================================
\begin{frame}{Expected Effects: Incidence of Non-Payment}
\vspace{-0.3cm}
\footnotesize

\blocklabel{Motivation}
Child support is \textbf{18\%} of single-parent families' resources, yet \textbf{30--40\%} of payments go unpaid: the loss falls on the mother or the State \parencite{bonnet2021flip}.

\begin{columns}[T,onlytextwidth]
  \begin{column}{0.41\textwidth}
    \blocklabel{Hypothesised mechanism}
    \begin{keybox}[Three-way incidence]
      \vspace{-0.3cm}
      \begin{align*}
        \text{Mother: } c^m &= z^m + p\,S + B\big(z^m,\tilde S\big) \\
        \text{Father: } c^f &= z^f - p\,S
      \end{align*}
      \vspace{-0.45cm}

      {\scriptsize $S$: support due; $p$: compliance, raised by default intermediation where the father can pay;
      $B$: ASF, RSA, prime d'activit\'e; $\tilde S$: support counted in means-testing; $\rho \equiv -\Delta B / (S\Delta p)$.}
      \vspace{-0.1cm}
      \[
        \underbrace{-\Delta c^f}_{\text{father pays}}
        = \underbrace{(1-\rho)\,S\,\Delta p}_{\text{mother gains}}
        + \underbrace{\rho\,S\,\Delta p}_{\text{State saves}}
      \]
      \vspace{-0.4cm}
    \end{keybox}
  \end{column}

  \begin{column}{0.56\textwidth}
    \blocklabel{Expected effects: heterogeneity across populations}
    \renewcommand{\arraystretch}{1.05}
    \begin{tabular}{@{}>{\raggedright\arraybackslash}p{0.36\linewidth}cc>{\raggedright\arraybackslash}p{0.40\linewidth}@{}}
      \toprule
      & $\Delta p$ & $\rho$ & Expected effect \\
      \midrule
      \multicolumn{4}{@{}l}{\color{boxframe}\textbf{Fathers}} \\
      Able to pay, defaulting & $> 0$ & & Pay more: $c^f \downarrow$ \\
      Unable to pay (low income) & $\approx 0$ & & No change; may cut formal earnings \parencite{cancian2013discouraging} \\
      \midrule
      \multicolumn{4}{@{}l}{\color{boxframe}\textbf{Mothers}} \\
      RSA recipients receiving support & & $\approx 1$ & Little change: the State pays less \\
      Above benefit thresholds & & $\approx 0$ & Net income gain \\
      Unpaid support counted ($\tilde S = S$) & & $\approx 0$ & Gain \\
      \bottomrule
    \end{tabular}

  \end{column}
\end{columns}
\end{frame}

% ============================================================
% SLIDE 5 -- Expected effects: labour supply
% ============================================================
\begin{frame}{Expected Effects: Labour Supply}
\vspace{-0.3cm}
\footnotesize

\blocklabel{Motivation}
Child support is a recurrent \emph{non-labour} income: making it larger and more reliable changes both the level and the volatility of resources.

\begin{columns}[T,onlytextwidth]
  \begin{column}{0.50\textwidth}
    \blocklabel{Hypothesised mechanisms}
    \begin{keybox}[Intensive margin: income effect]
      \vspace{-0.2cm}
      \begin{gather*}
        \max_{c,z}\; u(c,z) \quad \text{s.t.} \quad c = (1-\tau)\,z + R + (1-t_S)\,p\,S \\[0.1cm]
        \Delta z \simeq \frac{\eta}{1-\tau}\,(1-t_S)\,S\,\Delta p,
        \qquad \eta \equiv (1-\tau)\,\frac{\partial z}{\partial R} \le 0
      \end{gather*}
      \vspace{-0.35cm}

      {\scriptsize $t_S$: implicit tax on support ($t_S = 1$ for RSA recipients, $0$ above thresholds).}
    \end{keybox}
    \vspace{-0.2cm}
    \begin{keybox}[Extensive margin: stabilisation effect]
      \vspace{-0.2cm}
      \[
        \text{work} \iff u(c_1) - u(c_0) \geq q
      \]
      \vspace{-0.4cm}

      {\scriptsize $q$: fixed cost of work (e.g.\ childcare). Reliable support makes $q$ easier to cover $\Rightarrow$ more continuous employment.}
    \end{keybox}
  \end{column}

  \begin{column}{0.47\textwidth}
    \blocklabel{Expected effects}
    \begin{itemize}\setlength{\itemsep}{3pt}
      \item \textbf{Net effect ambiguous}: intensive $\downarrow$ vs.\ extensive $\uparrow$.
      \item \textbf{No income effect for RSA recipients} ($t_S = 1$): responses concentrated above benefit thresholds.
      \item \textbf{Per-euro effect} $\beta_z / \beta_S$, comparable with the MPE literature: \$1 of annual unearned income
            $\Rightarrow -\$0.52$ of labour earnings \textcite{golosov2024}.
            {\color{textgrey}Caveat: stabilisation works through regularity, not the amount (exclusion restriction).}
      \item \textbf{Gendered responses}: same incentive, different response by gender and gender norms
            \textcite{giommoni2025}.
    \end{itemize}
  \end{column}
\end{columns}
\end{frame}

% ============================================================
% SLIDE 6 -- Expected contributions
% ============================================================
\begin{frame}{Expected Contributions}
\vspace{-0.2cm}
\footnotesize
\begin{columns}[T,onlytextwidth]
\begin{column}{0.62\textwidth}
\begin{enumerate}
  \setlength{\itemsep}{0.1cm}
  \justifying
  \item \textbf{Causal evidence on default vs.\ opt-in enforcement.} We estimate the effect of making public
        intermediation the default, not of intermediation itself. The estimate is an ITT, which is policy-relevant
        because the default is the reform.
  \item \textbf{Incidence of non-payment.} We decompose mothers' post-separation resources into child support received
        and public transfers (ASF, RSA, prime d'activit\'e). This shows who bears unpaid support: the mother, the State or the father.
  \item \textbf{Labour supply channels.} We separate the income effect (intensive margin) from the stabilisation effect
        (extensive margin). Scaling the ITT by the first stage gives a per-euro effect that can be compared with the MPE
        literature \textcite{golosov2024}.
  \item \textbf{Design.} We use a difference-in-discontinuities with two cutoffs where treated and control groups swap
        roles, which allows each design to be checked against the other.
\end{enumerate}
\end{column}

\begin{column}{0.34\textwidth}
\begin{keybox}[What we do not do (scope)]
\begin{itemize}
  \setlength{\itemsep}{0.08cm}
  \item No effect of IFPA \emph{take-up} (individual IFPA status not observed)
  \item No outcomes for children
  \item No welfare analysis or optimal design of enforcement
  \item Short to medium run only ($\leq 24$ months)
\end{itemize}
\end{keybox}
\vspace{0.1cm}
\begin{keybox}[Who pays when fathers don't?]
  Likely mothers who bore defaults alone gain; for the poorest, the reform may only change \emph{who} pays.
\end{keybox}
\end{column}
\end{columns}
\end{frame}

% ============================================================
% REFERENCES (appendix, not counted in the slide numbers)
% ============================================================
\begin{frame}[noframenumbering]{References}
\vspace{-0.3cm}
\scriptsize
\begin{multicols}{2}
  \setlength{\parskip}{3pt}
  \printbibliography[heading=none]

  % Data sources (not in references.bib)
  \everypar{\hangindent=1.2em\hangafter=1}

  {\color{textgrey}Data sources (motivation): CNAF (2023); INSEE (2024); Minist\`ere de la Justice (2022).}
\end{multicols}
\end{frame}

\end{document}